% Source: https://artofproblemsolving.com/community/c3h1960661_prime_numbers
% Stable identity: aops-topic-1960661-q01
The prime digits are $2,3,5,7$. Rearranging,
\[
a(d+1)-2=d(15-b-2c). \tag{1}
\]
Reducing (1) modulo $d$ gives $d\mid a-2$. With $a,d\in\{2,3,5,7\}$, $a-2\in\{0,1,3,5\}$, so $a=2$ (any $d$), or $a=5,d=3$, or $a=7,d=5$.
\begin{itemize}
\item $a=2$: (1) reads $2d=d(15-b-2c)$, so $b+2c=13$. With prime digits, $b$ is odd: $(b,c)=(3,5)$ or $(7,3)$ ($b=5$ gives $c=4$). With $d$ arbitrary this gives $2352,2353,2355,2357,2732,2733,2735,2737$.
\item $a=5,d=3$: $18=3(15-b-2c)$, so $b+2c=9$: $(b,c)=(3,3)$ or $(5,2)$ ($b=7$ gives $c=1$). Numbers $5333$ and $5523$.
\item $a=7,d=5$: $40=5(15-b-2c)$, so $b+2c=7$: $(b,c)=(3,2)$ ($b=5$ gives $c=1$). Number $7325$.
\end{itemize}
So there are $11$ numbers:
\[
\boxed{2352,\ 2353,\ 2355,\ 2357,\ 2732,\ 2733,\ 2735,\ 2737,\ 5333,\ 5523,\ 7325}.
\]
Checked by computer.
